% TOP_PROBLEM: {"problemId": "problem.iwaniec-conjecture-for-the-beurling-ahlfors-transform-norm", "canonicalTitle": "Iwaniec Conjecture for the Beurling-Ahlfors Transform Norm", "releaseRank": 357, "primaryDomain": "analysis_pde", "primaryDomainLabel": "Analysis & PDE", "displayStatus": "open_with_solved_subcases", "releaseStatus": "open_with_solved_subcases"}
% !TeX program = lualatex
% ATTEMPT_STATUS: unresolved
% Model: GPT-6 Astra Ultra; continuation dated 27 September 2026.
\documentclass[11pt]{article}
\usepackage[margin=25mm]{geometry}
\usepackage{fontspec}
\setmainfont{Latin Modern Roman}
\usepackage{amsmath,amssymb,mathtools}
\usepackage{unicode-math}
\setmathfont{Latin Modern Math}
\usepackage{xurl,hyperref}
\hypersetup{colorlinks=true,urlcolor=blue}
\setcounter{secnumdepth}{0}
\setlength{\parindent}{0pt}
\setlength{\parskip}{5pt}
\setlength{\emergencystretch}{3em}
\begin{document}
\section*{\texorpdfstring{Iwaniec Conjecture for the Beurling-Ahlfors Transform Norm}{Iwaniec Conjecture for the Beurling-Ahlfors Transform Norm}}\label{357}
\subsection{Definitions and mathematical statement}
For suitable complex-valued functions on ${\mathbb{C}}\cong{\mathbb{R}}^2$, the Beurling--Ahlfors transform is the principal-value singular integral
\[
 Bf(z)=-\frac1\pi\operatorname{p.v.}\int_{{\mathbb{C}}}
 \frac{f(w)}{(z-w)^2}{\,\mathrm{d}} A(w).
\]
Its bounded extension to $L^p({\mathbb{C}})$, $1<p<\infty$, has norm $\|B\|_p=\sup_{f\ne0}\|Bf\|_p/\|f\|_p$. Put $p^*=\max\{p,p/(p-1)\}$. The conjecture is
\[
 \|B\|_p=p^*-1\quad\text{for every }1<p<\infty.
\]
The kernel normalization fixes the numerical constant. The functions are complex-valued and the norm is the full planar operator norm, not a restriction to radial or real-valued test functions. A nonsharp bound proportional to $p^*-1$ does not determine the requested value.
\par\smallskip\textit{Formulation and concept references:} \href{https://arxiv.org/html/2604.09365v2#S6}{[S1]}.

\subsection{Short English statement}
Is the exact norm of the Beurling--Ahlfors transform on every complex-valued $L^p$ space equal to $p^*-1$?
\subsection{Research attempt}
\paragraph{Scope and source audit.}
The normalization in the question is essential. We work with complex
functions, planar Lebesgue measure and the stated kernel. Section 6 of
the inspected 2026 version of [S1] still presents the sharp upper
bound as a problem. Its martingale results are not identified with a
proof of the conjecture. The elementary argument below proves the
entire proposed lower bound, proves equality at $p=2$, and reduces
the missing upper bound to a precise inequality for complex derivatives.
All these statements concern the full operator; restricting a test
family supplies a lower bound, not an upper bound on that operator.

Set
\[
 \partial=\tfrac12(\partial_x-i\partial_y),\qquad
 \bar\partial=\tfrac12(\partial_x+i\partial_y).
\]
Use the Fourier convention $\widehat f(\xi)=\int f(x)e^{-ix\cdot\xi}\,dx$.
The distributional fundamental solution
$\bar\partial(1/(\pi z))=\delta_0$ and its $\partial$ derivative
give the multiplier
\[
 m(\xi)=\frac{\xi_1-i\xi_2}{\xi_1+i\xi_2},\qquad \xi\ne0.
 \tag{357.1}
\]
For completeness, the fundamental-solution identity follows by
integrating against a smooth test function outside a circle of radius
$\varepsilon$ and applying the complex Green formula. The outer
boundary term vanishes for a compactly supported test function and
the inner boundary integral tends to its value at zero. Differentiating
$1/(\pi z)$ away from zero gives $-1/(\pi z^2)$, with the circular
principal value. Thus (357.1) uses exactly the kernel in the question.

\paragraph{The Hilbert-space case and the differential identity.}
Since $|m(\xi)|=1$ almost everywhere, Plancherel gives
\[
 \|Bf\|_2=\|f\|_2.
 \tag{357.2}
\]
In particular the sharp constant is one at $p=2$. This argument does
not extend to general $p$: bounded modulus of a Fourier multiplier
does not identify its $L^p$ norm unless $p=2$.

The same multiplier computation gives, for every
$u\in C_c^\infty(\mathbb C;\mathbb C)$,
\[
 B(\bar\partial u)=\partial u.
 \tag{357.3}
\]
There is no uncontrolled additive holomorphic function in this identity:
it is an equality of Fourier transforms, including equality as
distributions at zero frequency. Both derivatives are ordinary compactly
supported functions. Consequently any such $u$ gives the rigorous test
\[
 \|B\|_p\geq\frac{\|\partial u\|_p}{\|\bar\partial u\|_p}
 \quad\hbox{when }\bar\partial u\ne0.
 \tag{357.4}
\]

\paragraph{A cutoff power construction attaining the lower threshold.}
Fix $p\in(1,\infty)$ and put $\alpha=2/p\in(0,2)$.
First consider $v(z)=z|z|^{-\alpha}$ away from zero. Direct
differentiation, using $\partial |z|^2=\bar z$, gives
\[
 \partial v=(1-\alpha/2)|z|^{-\alpha},\qquad
 \bar\partial v=-\frac\alpha2\frac{z^2}{|z|^2}|z|^{-\alpha}.
 \tag{357.5}
\]
Neither derivative belongs to $L^p$ on the entire punctured plane:
the radial integral at the chosen exponent is logarithmic. We therefore
must cut off both ends rather than use $v$ as an inadmissible extremizer.

Choose smooth real radial cutoff functions $\chi_0,\chi_1$ such that
$\chi_0(r)=0$ for $r\leq1$, $\chi_0(r)=1$ for $r\geq2$,
$\chi_1(r)=1$ for $r\leq1$, and $\chi_1(r)=0$ for $r\geq2$.
For $0<\varepsilon<1/4$, define
\[
 u_\varepsilon(z)=\chi_0(|z|/\varepsilon)\chi_1(|z|)
                         z|z|^{-\alpha}.
\]
This is a smooth compactly supported complex function, identically zero
near the apparent singularity. On $2\varepsilon<|z|<1$ its derivatives
are exactly (357.5). Therefore the contributions from this annulus are
\[
 \begin{split}
 \int|\partial u_\varepsilon|^p
   &=2\pi(1-\alpha/2)^p\log(1/(2\varepsilon))+O_p(1),\\
 \int|\bar\partial u_\varepsilon|^p
   &=2\pi(\alpha/2)^p\log(1/(2\varepsilon))+O_p(1).
 \end{split}
 \tag{357.6}
\]
Here the bounded errors require verification. On the outer transition
$1\leq r\leq2$, every derivative is uniformly bounded. On the inner
transition $\varepsilon\leq r\leq2\varepsilon$, the product rule
bounds each derivative by $C_p\varepsilon^{-\alpha}$, while its area
is $O(\varepsilon^2)$. Its $p$-th power integral is bounded because
$\alpha p=2$. These are all remaining regions. Thus the errors in
(357.6) really are bounded independently of $\varepsilon$.

It follows from (357.4)--(357.6) that
\[
 \|B\|_p\geq\frac{1-\alpha/2}{\alpha/2}=p-1.
 \tag{357.7}
\]
Apply the same construction with $\bar z|z|^{-\alpha}$, or simply
conjugate $u_\varepsilon$. Complex conjugation exchanges the moduli
of $\partial u$ and $\bar\partial u$. It therefore gives the second
lower bound
\[
 \|B\|_p\geq\frac1{p-1}.
 \tag{357.8}
\]
Taking the larger of (357.7) and (357.8) proves
\[
 \boxed{\ \|B\|_p\geq p^*-1\quad(1<p<\infty).\ }
 \tag{357.9}
\]
This proof also explains the change of the relevant test at $p=2$:
one orientation of the complex power is useful above two, and its
conjugate is useful below two. A calculation with only one orientation
would miss half of the proposed lower bound.

\paragraph{Duality and why it does not provide the upper bound.}
Write $Cf=\bar f$. Since $m(-\xi)=m(\xi)$, Fourier transformation
shows $CBC=B^*$ on $L^2$, where the latter has multiplier $\bar m$.
These identities extend to the bounded $L^p$ realizations by density.
Conjugation is an isometry. If $q=p/(p-1)$, Banach-space duality gives
\[
 \|B\|_p=\|B^*\|_q=\|B\|_q.
 \tag{357.10}
\]
Thus proving the conjectured upper bound for $p\geq2$ would settle
the remaining range. This equality is a reduction, not an estimate:
it only relates two currently unknown norms. The lower-bound tests
already respect it.

There is also a useful identity at the level of Jacobians. If
$u=u_1+iu_2$, direct expansion gives
\[
 |\partial u|^2-|\bar\partial u|^2
   =\partial_xu_1\partial_yu_2-\partial_yu_1\partial_xu_2.
 \tag{357.11}
\]
For a compactly supported smooth $u$ the integral of the right-hand
side is zero by integration by parts, recovering equality of the two
$L^2$ derivative norms. A tempting argument multiplies (357.11) by
a power of $|Du|$ and tries the same integration. It fails because the
new weight has nonzero derivatives. There is no longer a null
Lagrangian with zero integral. The logarithmic tests make clear that
any replacement must permit derivative-norm ratios as large as
$p^*-1$.

\paragraph{An exact analytic form of the missing estimate.}
For fixed $p$, the conjectured upper bound would imply
\[
 \int_{\mathbb C}|\partial u|^p\,dA
   \leq(p^*-1)^p\int_{\mathbb C}|\bar\partial u|^p\,dA
 \quad(u\in C_c^\infty(\mathbb C;\mathbb C)).
 \tag{357.12}
\]
Conversely, this family of derivative inequalities is sufficient for
the full operator upper bound. Here is the density point needed in
that assertion. The subspace
$\{\bar\partial u:u\in C_c^\infty\}$ is dense in $L^p$.
Indeed, if a continuous dual functional annihilates it, its representing
$L^q$ function solves one of the complex Cauchy--Riemann equations in
distributions. Mollification gives an entire holomorphic or
antiholomorphic function with bounded $L^q$ norm. The mean-value
inequality on a disc of radius $R$ bounds its value at any fixed point
by $C R^{-2/q}$ times that norm. Letting $R$ tend to infinity makes
the mollified function zero; then mollifiers tending to the identity
make the original dual function zero. Hahn--Banach proves density.
The possible conjugation in the dual pairing only exchanges the two
Cauchy--Riemann equations and does not change this argument.

Now apply (357.12) and (357.3) on this dense subspace and use the
already known bounded extension of $B$ on $L^p$. Passage to the
$L^p$ limit yields the desired bound for every input. Thus (357.12)
is a precise first missing statement; it is not enough to prove it
for radial $u$ or for a single angular mode.

\paragraph{Why mode computations and discretization do not finish it.}
Rotations decompose $L^2$ orthogonally into angular Fourier modes,
and (357.1) shifts the angular behavior in a particularly simple way.
For $p\ne2$, however,
\[
 \|f_1+f_2\|_p^p\ne\|f_1\|_p^p+\|f_2\|_p^p
\]
in general, even for different angular frequencies. Bounds proved on
each mode therefore cannot be summed with constant one. Cancellations
in the input and reinforcement in the output are precisely what a
full-space upper estimate must control. An operator matrix truncated
in angular and radial frequency likewise has a restricted test space;
its computed norm is not an upper bound on the infinite-dimensional
operator without an independent uniform remainder estimate.

The associated diagnostic checks the derived power-coefficient ratios
with exact rational arithmetic for selected $p$, and checks
$|m(\xi)|=1$ numerically on nonzero integer frequencies. It does not approximate
a proof of (357.12). The annular cutoff estimate is analytic and
uniform, so the lower bound is not inferred from numerical fitting.

\paragraph{Outcome.}
The lower threshold is rigorously necessary at every $p$, and it is
the exact norm at $p=2$. Duality reduces the open upper estimate to
one side of two, and density identifies its full differential form.
The first gap is establishing (357.12) for arbitrary complex compactly
supported maps. Neither the $L^2$ cancellation, the power tests nor
an angular-mode restriction supplies that estimate. The full claim
remains unresolved by this attempt.
\subsection{Sources}
\begin{itemize}
\item[S1]Rodrigo Banuelos, Cotlar martingale transforms and related singular integrals, arXiv:2604.09365v2 (2026).
\url{https://arxiv.org/html/2604.09365v2#S6}.
\textit{Formulation source.}
\end{itemize}
\end{document}
